% GENERATED FILE. Edit canonical lessons, then run npm run generate:books.
% canonical-source-hash: fnv1a64:a58953cd17be27d0
% canonical-lessons: SA-C206-majjati, SA-C206-snati, SA-C206-khelati, SA-C206-nimilayati, SA-C206-unmilayati

\chapter{Sinks, Bathes, Frolics, Shuts, Opens}
\label{ch:sa-sinks-bathes-frolics-shuts-opens}
\hlchaptermodality{206}

\begin{chapteropening}
\textbf{By the end of this chapter:} \emph{I can say sinks, bathes, frolics, shuts and opens.}

Complete the last lesson of chapter 206: I can say sinks, bathes, frolics, shuts and opens.
\end{chapteropening}

\section[\texorpdfstring{\sk{मज्जति} (majjati)}{मज्जति (majjati)}]{\sk{मज्जति} (majjati) — sinks, dives}

\label{lesson:SA-C206-majjati}



\begin{quote}
Before the new one: say the Sanskrit for accepts, then the Sanskrit for blossoms.
\end{quote}

\subsection*{You'll want to know: \sk{मज्जति}}
\textbf{\sk{मज्जति}} — \emph{majjati} — \textquotedblleft{}sinks, dives\textquotedblright{}.

\begin{cousinweb}
\emph{majjati} is the present of the root √majj, \textquotedblleft{}to sink, dive\textquotedblright{} (MW). In the Rigveda it means \textquotedblleft{}sinks, perishes\textquotedblright{} (CDIAL 9711). MW compares Latin \emph{mergere}, \textquotedblleft{}to dip, sink\textquotedblright{}, the source of English \emph{merge} and \emph{submerge}.
\end{cousinweb}

\begin{grammarlens}[title={What you've built}]
One more word for this chapter.
\end{grammarlens}

\subsection*{Your turn}
\begin{itemize}
\raggedright
  \item \emph{Say it:} \emph{majjati}
  \item \emph{Say it:} \emph{majjati}, once more
  \item \emph{Say it:} say \emph{majjati}
  \item \emph{Recall:} say the Sanskrit for accepts, then the Sanskrit for blossoms, then say \emph{majjati} again
\end{itemize}

\begin{tcolorbox}[breakable,colback=teal!4,colframe=teal!35!black,title={Before you move on}]
What does \sk{मज्जति} mean? (\textquotedblleft{}Sinks, dives\textquotedblright{}.) Say it once more.
\end{tcolorbox}
\section[\texorpdfstring{\sk{स्नाति} (snāti)}{स्नाति (snāti)}]{\sk{स्नाति} (snāti) — bathes}

\label{lesson:SA-C206-snati}



\begin{quote}
Before the new one: say the Sanskrit for blossoms, then the Sanskrit for sinks.
\end{quote}

\subsection*{You'll want to know: \sk{स्नाति}}
\textbf{\sk{स्नाति}} — \emph{snāti} — \textquotedblleft{}bathes\textquotedblright{}.

\begin{cousinweb}
\emph{snāti} is the present of the root √snā, \textquotedblleft{}to bathe\textquotedblright{} (MW). The root is inherited from Indo-European: MW compares Latin \emph{nāre}, \textquotedblleft{}to swim\textquotedblright{}. It lives on as Hindi \emph{nahānā}, \textquotedblleft{}to bathe\textquotedblright{} (CDIAL 13786).
\end{cousinweb}

\begin{grammarlens}[title={What you've built}]
One more word for this chapter.
\end{grammarlens}

\subsection*{Your turn}
\begin{itemize}
\raggedright
  \item \emph{Say it:} \emph{snāti}
  \item \emph{Say it:} \emph{snāti}, once more
  \item \emph{Say it:} say \emph{snāti}
  \item \emph{Recall:} say the Sanskrit for blossoms, then the Sanskrit for sinks, then say \emph{snāti} again
\end{itemize}

\begin{tcolorbox}[breakable,colback=teal!4,colframe=teal!35!black,title={Before you move on}]
What does \sk{स्नाति} mean? (\textquotedblleft{}Bathes\textquotedblright{}.) Say it once more.
\end{tcolorbox}
\section[\texorpdfstring{\sk{खेलति} (khelati)}{खेलति (khelati)}]{\sk{खेलति} (khelati) — frolics, plays}

\label{lesson:SA-C206-khelati}



\begin{quote}
Before the new one: say the Sanskrit for sinks, then the Sanskrit for bathes.
\end{quote}

\subsection*{You'll want to know: \sk{खेलति}}
\textbf{\sk{खेलति}} — \emph{khelati} — \textquotedblleft{}frolics, plays\textquotedblright{}.

\begin{cousinweb}
\emph{khelati} first means \textquotedblleft{}shakes, moves to and fro\textquotedblright{} (MW), and later \textquotedblleft{}plays\textquotedblright{}. CDIAL 3918 says it may go with a family of \textquotedblleft{}play\textquotedblright{} words that are perhaps non-Aryan; their origin is uncertain. Hindi \emph{khelnā}, \textquotedblleft{}to play\textquotedblright{}, belongs to that family.
\end{cousinweb}

\begin{grammarlens}[title={What you've built}]
One more word for this chapter.
\end{grammarlens}

\subsection*{Your turn}
\begin{itemize}
\raggedright
  \item \emph{Say it:} \emph{khelati}
  \item \emph{Say it:} \emph{khelati}, once more
  \item \emph{Say it:} say \emph{khelati}
  \item \emph{Recall:} say the Sanskrit for sinks, then the Sanskrit for bathes, then say \emph{khelati} again
\end{itemize}

\begin{tcolorbox}[breakable,colback=teal!4,colframe=teal!35!black,title={Before you move on}]
What does \sk{खेलति} mean? (\textquotedblleft{}Frolics, plays\textquotedblright{}.) Say it once more.
\end{tcolorbox}
\section[\texorpdfstring{\sk{निमीलयति} (nimīlayati)}{निमीलयति (nimīlayati)}]{\sk{निमीलयति} (nimīlayati) — shuts (the eyes)}

\label{lesson:SA-C206-nimilayati}



\begin{quote}
Before the new one: say the Sanskrit for bathes, then the Sanskrit for frolics.
\end{quote}

\subsection*{You'll want to know: \sk{निमीलयति}}
\textbf{\sk{निमीलयति}} — \emph{nimīlayati} — \textquotedblleft{}shuts (the eyes)\textquotedblright{}.

\begin{cousinweb}
\emph{nimīlayati} is built on \emph{ni-}, \textquotedblleft{}down\textquotedblright{}, and the causative of the root √mīl, \textquotedblleft{}to close the eyes\textquotedblright{} (MW; CDIAL 10143). The Rigveda already has \emph{sam mīlya}, \textquotedblleft{}with closed eyes\textquotedblright{}. MW gives the root no older source.
\end{cousinweb}

\begin{grammarlens}[title={What you've built}]
One more word for this chapter.
\end{grammarlens}

\subsection*{Your turn}
\begin{itemize}
\raggedright
  \item \emph{Say it:} \emph{nimīlayati}
  \item \emph{Say it:} \emph{nimīlayati}, once more
  \item \emph{Say it:} say \emph{nimīlayati}
  \item \emph{Recall:} say the Sanskrit for bathes, then the Sanskrit for frolics, then say \emph{nimīlayati} again
\end{itemize}

\begin{tcolorbox}[breakable,colback=teal!4,colframe=teal!35!black,title={Before you move on}]
What does \sk{निमीलयति} mean? (\textquotedblleft{}Shuts (the eyes)\textquotedblright{}.) Say it once more.
\end{tcolorbox}
\section[\texorpdfstring{\sk{उन्मीलयति} (unmīlayati)}{उन्मीलयति (unmīlayati)}]{\sk{उन्मीलयति} (unmīlayati) — opens (the eyes)}

\label{lesson:SA-C206-unmilayati}



\begin{quote}
Before the new one: say the Sanskrit for frolics, then the Sanskrit for shuts.
\end{quote}

\subsection*{You'll want to know: \sk{उन्मीलयति}}
\textbf{\sk{उन्मीलयति}} — \emph{unmīlayati} — \textquotedblleft{}opens (the eyes)\textquotedblright{}.

\begin{cousinweb}
\emph{unmīlayati} is built on \emph{ud}, \textquotedblleft{}up\textquotedblright{}, and the causative of the root √mīl, \textquotedblleft{}to close the eyes\textquotedblright{} (MW). \emph{Ud} turns closing into opening, the mirror image of \emph{nimīlayati}.
\end{cousinweb}

\begin{grammarlens}[title={What you've built}]
One more word for this chapter.
\end{grammarlens}

\subsection*{Your turn}
\begin{itemize}
\raggedright
  \item \emph{Say it:} \emph{unmīlayati}
  \item \emph{Say it:} \emph{unmīlayati}, once more
  \item \emph{Say it:} say \emph{unmīlayati}
  \item \emph{Recall:} say the Sanskrit for frolics, then the Sanskrit for shuts, then say \emph{unmīlayati} again
\end{itemize}

\begin{tcolorbox}[breakable,colback=teal!4,colframe=teal!35!black,title={Before you move on}]
What does \sk{उन्मीलयति} mean? (\textquotedblleft{}Opens (the eyes)\textquotedblright{}.) Say it once more.
\end{tcolorbox}
